\documentclass[10pt,twocolumn]{article}
\usepackage[margin=0.6in]{geometry}
\usepackage[T1]{fontenc}
\usepackage[utf8]{inputenc}
\usepackage{lmodern}
\usepackage{amsmath}
\usepackage{graphicx}
\usepackage{booktabs}
\usepackage{array}
\usepackage{tabularx}
\usepackage[table]{xcolor}
\usepackage{enumitem}
\usepackage{caption}
\usepackage{tikz}
\usepackage{float}
\usepackage{placeins}
\usepackage{url}
\usepackage[hidelinks]{hyperref}
\usetikzlibrary{arrows.meta,positioning,calc}

\captionsetup{font=footnotesize,labelfont=bf}
\setlength{\columnsep}{0.20in}
\setlength{\parindent}{0pt}
\setlength{\parskip}{0.32em}
\raggedbottom

\title{\vspace{-1.5em}\textbf{Question Answering with a Custom Transformer Encoder and Generative Decoder}\vspace{-0.5em}}
\author{\href{https://github.com/Harsha081459/Question-Answer-System-with-a-Custom-Transformer-Encoder-and-Generative-Decoder}{\textbf{[GitHub]}} \quad $\mid$ \quad \href{https://iiitbac-my.sharepoint.com/:f:/g/personal/vishal_sriram_iiitb_ac_in/IgBJMBmHEodfQba6fYpup9aiAdhsUH577v5oWa_W0i20VqM?e=j7DCc6}
{\textbf{[Weights -- OneDrive]}} \quad $\mid$ \quad \href{https://huggingface.co/spaces/hv-123/QA-Engine}{\textbf{[Web App]}}}

\date{}

\begin{document}
\maketitle
\vspace{-2.5em}

% ── Team Information ──
\begin{table}[H]
\centering
\small
\caption*{\textbf{NLP (AID 829) --- Group Assignment $\mid$ Group 18}}
\vspace{-0.5em}
\begin{tabular}{@{}clc@{}}
\toprule
\textbf{S.No.} & \textbf{Name} & \textbf{Roll Number} \\
\midrule
1 & Vishal Sriram & BT2024036 \\
2 & Harsha & BT2024106 \\
3 & Lohith & BT2024248 \\
4 & Munjikesh & BT2024247 \\
\bottomrule
\end{tabular}
\end{table}
\vspace{-0.5em}

% ══════════════════════════════════════════════════════════════
\begin{abstract}
We built a SQuAD-based question answering system featuring two parallel tracks: an extractive span predictor and a custom generative answer decoder. Both tracks share a custom 12-layer Transformer encoder pretrained from scratch via masked language modeling. The extractive model was benchmarked against strong public baselines, confirming that our scratch-built encoder learns a competitive span-selection representation. For the generative track, we developed a bespoke 4-layer decoder optimized for concise answer generation. Evaluated on a shared validation subset, our custom decoder outperformed T5-small and Flan-T5-small, closely approaching T5-base.
\end{abstract}

% ══════════════════════════════════════════════════════════════
\section{Introduction}
Machine reading comprehension requires a model to process a question and context passage, then return a precise answer. This project addressed the task through two distinct modes: \textbf{extractive QA} (predicting start and end token indices of the answer span) and \textbf{generative QA} (decoding a short answer token-by-token from the encoder states). Implementing both modes provided a rigorous test of the underlying encoder's representational capacity.

The pipeline executed in four sequential stages:
\begin{enumerate}[leftmargin=1.2em,itemsep=3pt,topsep=4pt]
    \item \textbf{Stage 1:} Pretrain a BERT-style encoder from scratch using masked language modeling on Wikipedia + C4.
    \item \textbf{Stage 2:} Fine-tune the encoder for extractive QA on SQuAD v2.
    \item \textbf{Stage 3:} Construct a generative encoder-decoder on top of the same encoder.
    \item \textbf{Stage 4:} Expose both inference paths through a unified local demo.
\end{enumerate}

\vfill\columnbreak

\textbf{Primary contribution.} A clear engineering and research progression: a single custom-trained encoder powering both accurate extractive span prediction and competitive generative answer decoding without relying on off-the-shelf pretrained architectures.

% ══════════════════════════════════════════════════════════════
\section{Related Work}
Extractive QA was revolutionized by BERT \cite{devlin2019bert}, which introduced bidirectional Transformer pretraining via masked language modeling. Our encoder architecture closely follows the BERT-base specification: a 12-layer Transformer stack with 768-dimensional hidden states. While robustly optimized variants like RoBERTa \cite{liu2019roberta} achieve higher absolute performance on SQuAD \cite{rajpurkar2016squad}, our project focused on the pedagogical value of pretraining this architecture from scratch.

In generative QA, unified text-to-text Transformers like T5 \cite{raffel2020t5} dominate. We use T5-small and T5-base as external benchmarks, but our approach differs by grafting a lighter, bespoke 4-layer decoder onto our BERT-style encoder via a linear projection bridge. This hybrid design allowed strict control over decoding heuristics without inheriting the broad generation priors of massively pretrained seq2seq models.

% ══════════════════════════════════════════════════════════════
\section{System Architecture}

The system was designed around a single shared encoder whose contextual
representations power two structurally distinct inference heads. This
unified backbone architecture reflects a deliberate engineering choice:
rather than training separate models for extractive and generative QA,
we exploit the fact that both tasks require the same fundamental
capability --- grounding token-level representations in the joint
context of a question and a supporting passage.

\subsection{Input Representation}
Both inference paths begin with identical preprocessing. The question
and context are concatenated using the standard BERT-style
\texttt{[CLS] Q [SEP] C [SEP]} template. Three embedding matrices are
summed to form the initial token representation: a token embedding
($|\mathcal{V}| \times 768$), a learned positional embedding (up to
512 positions), and a segment embedding that distinguishes question
tokens (type~0) from context tokens (type~1). The summed embeddings
are passed through layer normalization and dropout before entering the
encoder stack.

\subsection{Shared Transformer Encoder}
The encoder is a 12-layer bidirectional Transformer. Each layer
applies pre-layer-norm multi-head self-attention (12 heads,
$d_k = d_v = 64$) followed by a two-layer feed-forward network with
hidden dimension 3072 and GELU activation. Residual connections wrap
both sub-layers. The full sequence output
$\mathbf{H} \in \mathbb{R}^{L \times 768}$ is passed to both
downstream heads, where $L$ is the padded sequence length (capped
at 256 tokens with a doc stride of 64 for long contexts).

The pre-norm variant was chosen over the original post-norm BERT
design because it yields more stable gradients during fine-tuning,
particularly when the encoder is jointly optimized with the decoder
at a lower learning rate.

\subsection{Extractive Head}
The extractive head is a single linear projection
$\mathbf{W} \in \mathbb{R}^{768 \times 2}$ applied independently to
every encoder output token. This produces two scalar logits per
position: $s_i$ (start score) and $e_i$ (end score). At inference
time, the top-$n_{\text{best}}$ start positions are enumerated and
paired with all legal end positions to form candidate spans. Each span is scored by summing its start and end logits, and the highest-scoring valid
span is returned as the extracted answer. For unanswerable questions
(SQuAD v2 no-answer instances), the score of the \texttt{[CLS]} token
logit serves as the null-answer score; if it exceeds the best span
score by a tuned threshold $\tau$, the model abstains.

\subsection{Encoder--Decoder Bridge}
Because the encoder operates in a 768-dimensional space while the
decoder uses a 512-dimensional space, a learned linear projection
$\mathbf{W}_{\text{bridge}} \in \mathbb{R}^{768 \times 512}$ maps
the full encoder memory $\mathbf{H}$ into the decoder's key--value
space before each cross-attention layer. This projection is trained
jointly with the decoder and allows the two modules to operate at
different widths without any information bottleneck beyond what
the linear map can express.

\subsection{Generative Decoder}
The generative decoder is a 4-layer autoregressive Transformer with
hidden dimension 512 and 8 attention heads ($d_k = d_v = 64$). Each
decoder block applies three sub-layers in sequence, each wrapped in a
pre-norm residual:

\begin{enumerate}[leftmargin=1.2em,noitemsep,topsep=2pt]
    \item \textbf{Masked self-attention.} Attends over all previously
    generated tokens using a causal mask, building up a
    representation of the partial answer prefix.
    \item \textbf{Cross-attention.} Queries the projected encoder
    memory $\mathbf{W}_{\text{bridge}}\mathbf{H}$ to ground each
    decoding step in the source context, directly suppressing
    hallucination.
    \item \textbf{Feed-forward network.} A two-layer FFN (hidden dim
    2048, GELU) refines the attended representation before the


# (SQuAD v2 no-answer instances), the score of the \texttt{[CLS]} token
# logit serves as the null-answer score; if it exceeds the best span
# score by a tuned threshold $\tau$, the model abstains.
# 
# \subsection{Encoder--Decoder Bridge}
# Because the encoder operates in a 768-dimensional space while the
# decoder uses a 512-dimensional space, a learned linear projection
# $\mathbf{W}_{\text{bridge}} \in \mathbb{R}^{768 \times 512}$ maps
# the full encoder memory $\mathbf{H}$ into the decoder's key--value
# space before each cross-attention layer. This projection is trained
# jointly with the decoder and allows the two modules to operate at
# different widths without any information bottleneck beyond what
# the linear map can express.
# 
# \subsection{Generative Decoder}
# The generative decoder is a 4-layer autoregressive Transformer with
# hidden dimension 512 and 8 attention heads ($d_k = d_v = 64$). Each
# decoder block applies three sub-layers in sequence, each wrapped in a
# pre-norm residual:
# 
# \begin{enumerate}[leftmargin=1.2em,noitemsep,topsep=2pt]
#     \item \textbf{Masked self-attention.} Attends over all previously
#     generated tokens using a causal mask, building up a
#     representation of the partial answer prefix.
#     \item \textbf{Cross-attention.} Queries the projected encoder
#     memory $\mathbf{W}_{\text{bridge}}\mathbf{H}$ to ground each
#     decoding step in the source context, directly suppressing
#     hallucination.
#     \item \textbf{Feed-forward network.} A two-layer FFN (hidden dim
#     2048, GELU) refines the attended representation before the
